\section{Discussion}

The preceding sections introduced the AnyBlox architecture and its decoding model. Whether this design is practical depends on two questions:
\begin{enumerate}
    \item What overhead does the WebAssembly sandbox introduce?
    \item Where does the approach reach its limits?
\end{enumerate}
Gienieczko et al. (2025) report benchmarks that cover both favorable and unfavorable cases for AnyBlox.

\subsection{Performance Evaluation}

Sandboxing a decoder inside Wasm adds a layer of indirection between the data and the query engine. Gienieczko et al. (2025) measure the resulting overhead directly. They run TPC-H at scale factor 20 in DataFusion~\cite{lambApacheArrowDataFusion2024} and vary the representation of \texttt{lineitem}, the largest table in the benchmark and therefore the one whose scan dominates the runtime, across three configurations:
\begin{description}
    \item[Parquet] read by DataFusion's native reader with Snappy\footnote{Snappy, \url{https://github.com/google/snappy/} (accessed: June 15, 2026)} compression;
    \item[AnyBlox Vortex] a Vortex file decoded inside the AnyBlox sandbox;
    \item[Native Vortex] the same Vortex decoder compiled natively without any sandbox.
\end{description}
The first variant isolates the effect of the format; the latter two isolate the effect of the sandbox alone. Predicate pushdown was disabled throughout, so the numbers reflect scan and operator cost only. Vortex\footnote{Vortex, \url{https://vortex.dev/} (accessed: June 15, 2026)} is a recent columnar format.

The benchmark yields two results. Both Vortex variants achieve over $2\times$ lower scan latency than Parquet while also reaching a better compression rate. For the Parquet reader, decoding accounts for nearly half of the processing time, while Vortex spends around 80\% of its processing time in relational operators. The overhead of the sandbox itself is small. Sandboxed Vortex decoding adds 5\% to the scan latency over native decoding, which amounts to a 1\% increase in overall query latency and which the authors report as within measurement noise. Across this end-to-end workload, the cost of the sandbox is therefore negligible, and AnyBlox provides DataFusion with an encoding that outperforms its own native Parquet reader~\cite{gienieczkoAnyBloxFrameworkSelfDecoding2025}.